% 自定义彩色圆圈数字命令
\newcommand{\circled}[1]{\textcolor{blue}{\textcircled{\textcolor{black}{\small #1}}}}
%\newcommand{\circled}[1]{\textcircled{\small #1}}

% 重新定义 notes 环境
\newenvironment{notes}{
	\par\noindent\makebox[0pt][r]{%
		\scriptsize\color{red!90}\textdbend\quad}% 图标
	\textbf{\color{red!90}Note:} \normalfont % 红色提示标题，正体字体
}{\par}

% 定义 change 环境，使用圆圈数字标注每一个变化项
%\newenvironment{change}{
%	\begin{enumerate}[label=\small\protect\circled{\arabic*}]}{
%\end{enumerate}}

\newenvironment{change}{
	\begin{enumerate}[label=\hfill\circled{\arabic*}, itemsep=1em]
	}{
	\end{enumerate}
}
% 定义 change2 环境
\newenvironment{change2}{
	\begin{tcolorbox}[colback=white, colframe=gray, boxrule=0.5mm, sharp corners=south]
		\begin{enumerate}[label=\centering\circled{\arabic*}]
		}{
		\end{enumerate}
	\end{tcolorbox}
}
\newenvironment{kuohao}{
	\begin{enumerate}[label=\textbf{(\arabic*)}, itemsep=1em]  % Using \textbf to make the numbers bold and enclosed in parentheses
	}{
	\end{enumerate}
}
% 定义 tip 环境，用于显示提示信息
\newenvironment{tip}{
	\par\noindent\makebox[-3pt][r]{%
		\scriptsize\color{green!90}\ding{43}\quad}% 绿色图标
	\textbf{\color{green!90}Tip:} \itshape % 绿色的提示标题
}{\par}

% 定义 warning 环境，用于显示警告信息
\newenvironment{warning}{
	\par\noindent\makebox[-3pt][r]{%
		\scriptsize\color{orange!90}\ding{70}\quad}% 橙色图标
	\textbf{\color{orange!90}Warning:} \itshape % 橙色的警告标题
}{\par}

% 定义 important 环境，用于显示重要信息
\newenvironment{important}{
	\par\noindent\makebox[-3pt][r]{%
		\scriptsize\color{red!90}\ding{72}\quad}% 红色图标
	\textbf{\color{red!90}Important:} \itshape % 红色的重点标题
}{\par}

% 定义 examplezero 环境，并提供自动编号
\newcounter{examplezero}[section]
\renewcommand{\theexamplezero}{\thesection.\arabic{examplezero}}

\newenvironment{examplezero}[1][]{
	\refstepcounter{examplezero} % 增加 examplezero 计数器
	\par\noindent\textbf{\color{purple!90}Example \theexamplezero:} \itshape #1 \rmfamily}{\par}
	
% 定义introduction环境
\newenvironment{introduction}[1][Introduction]{
	\begin{tcolorbox}[title={#1}]
		\begin{multicols}{2}
			\begin{itemize}[label=\textcolor{structurecolor}{\upshape\scriptsize$\bullet$}]  % 使用$\bullet$
			}{
			\end{itemize}
		\end{multicols}
	\end{tcolorbox}
}

% 上标文献
\newcommand{\upcite}[1]{$^{\mbox{\scriptsize \cite{#1}}}$}
\newcommand{\supercite}[1]{\textsuperscript{\textsuperscript{\cite{#1}}}}
%------------左边放图，右边说明----
\newenvironment{leftimage}[2][0.45\textwidth]{% #1: 图像宽度，#2: 文字说明
	\begin{figure}[htbp]
		\centering
		\begin{minipage}{#1}  % 左边的图像区域
			\centering
			\includegraphics[width=\linewidth]{#2}  % 图片路径
		\end{minipage}%
		\hspace{0.05\textwidth}  % 左右两部分间距
		\begin{minipage}{0.45\textwidth}  % 右边的文字说明区域
			\raggedright
		}{%
		\end{minipage}
	\end{figure}
}
%%%---
% 自定义环境：左图右文
\newenvironment{lefttwoimage}[2][0.45\textwidth]{% #1: 图像宽度，#2: 图片名称
	\begin{figure}[htbp]
		\centering
		\begin{minipage}{#1}  % 左边的图像区域
			\centering
			\includegraphics[width=\linewidth]{#2}  % 图片路径
		\end{minipage}%
		\hspace{0.05\textwidth}  % 左右两部分间距
		\begin{minipage}{0.45\textwidth}  % 右边的文字说明区域
			\raggedright
		}{%
		\end{minipage}
	\end{figure}
}
%--% 自定义环境：左图右文，支持标签和引用
% 自定义环境：左图右文，支持标签和引用
\newenvironment{leftthreeimage}[2][0.45\textwidth]{% #1: 图像宽度，#2: 图片名称
	\begin{figure}[htbp]
		\centering
		\begin{minipage}{#1}  % 左边的图像区域
			\centering
			\includegraphics[width=\linewidth]{#2}  % 图片路径
		\end{minipage}%
		\hspace{0.05\textwidth}  % 左右两部分间距
		\begin{minipage}{0.45\textwidth}  % 右边的文字说明区域
			\raggedright
		}{%
		\end{minipage}
		% 用户可以在此自定义标题和标签
	\end{figure}
}

